\documentclass[~/PhyAI/phyai.tex]{subfiles}
\begin{document}
\chapter{Magnetism in Matter}\label{chap:magnetism_of_matters} % (fold)
\epigraph{I think I can safely say that nobody understands quantum mechanics.}{\textit{Richard Feynman}}
\section{Magnetic Dipoles and Diamagnetism}\label{sec:magnetic_dipoles_and_diamagnetic} % (fold)
\subsection{Magnetic Dipoles}
I think this final chapter is the most subtle in the book. From the first chapter until now, we have only been working with the assumption that all matter is made up of randomly moving "particles" without considering their structure. Do you remember Atomic Theory from the previous chapter (Theorem ~\ref{Atomic theory})? We treat it as the foundation (or as I said, an axiom). But we are having trouble explaining the essence of magnetism, our old theory, simply reaches its own limit. Now we must use a bit of the Bohr's model to keep everything on track. Let's begin with a simple experiment, imagine you have a single loop of wire between two magnetic poles. Inside the loop, a current flows counterclockwise. Applying left-hand rule, you can draw the direction of Lorentz force, which causes the loop to rotate:
\begin{figure}[H]
	\begin{center}
		\includegraphics[width=0.5\textwidth]{~/PhyAI/chapter6/6_1.png}
	\end{center}
	\caption{A loop inside magnetic field $\vec{B}$}\label{fig:6_1}
\end{figure}
A rectangular loop has size $a\times b$, under the net torque:
$$ \vec{\tau}_{\text{net}}=\sum\vec{\tau}=\sum\vec{R}\times\vec{F} =\sum\vec{R}\times i\dot(\vec{L}\times\vec{B})$$
I define $\vec{n}$ as the normal vector of the loop, and an angle $\theta=(\vec{n},\vec{B})$, because both torques vectors are in the same direction, we obtain:
\begin{figure}[H]
	\begin{center}
		\includegraphics[width=0.5\textwidth]{~/PhyAI/chapter6/6_2.png}
	\end{center}
	\caption{Both torque vectors in pointing out from the plane of paper}\label{fig:6_2}
\end{figure}
The scalar $\tau_{net}$ value is:
$$ \tau_{\text{net}}=2.\frac{b }{2}.iaB\sin\theta=iab\sin\theta.B=iA\sin\theta.B $$
I define a new vector $\vec{\mu}$ quantity, called \textbf{magnetic dipole moment}, same direction to normal vector $\vec{n}$, and has magnitude: $$ |\vec{\mu}|=iA =iab$$
Rewriting the net torque equation yields:
$$ \vec{\tau}_{\text{net}}=\vec{\mu}\times \vec{B} $$
\begin{definition}[Magnetic Moment Definition]
	The \textbf{magnetic moment} (or magnetic dipole moment) is a vector quantity that determines the magnitude of torque in a given magnetic field:
	$$ \vec{\tau}=\vec{\mu}\times  \vec{B}$$
\end{definition}
Each angle $\theta$ (position), it is characterized by a \textbf{potential energy value}. I randomly choose 2 angles $\theta_{i}$ and $\theta_{f}$:
$$ W=-\Delta U=\int_{\theta_{i}}^{\theta_{f}}\tau d\theta=-|\vec{\mu}|.|\vec{B}|.(\cos\theta_{f}-\cos\theta_{i})=-(U_{f}-U_{i})$$
I assign the potential energy value to angle $\theta$: $$ U(\theta)=-\vec{\mu}\cdot\vec{B} $$
From here we have a very important result: a magnetic dipole has its lowest energy when its vector $\vec{\mu}$ is lined up with magnetic field $\vec{B}$, and vice versa.
\begin{figure}[H]
	\begin{center}
		\includegraphics[width=0.666666\textwidth]{~/PhyAI/chapter6/6_3.png}
	\end{center}
	\caption{Highest and lowest potential energy value}\label{fig:6_3}
\end{figure}
To flip an arrow (magnetic moment vector) down, we must supply $2\mu B$ external energy units.
% section Magnetic Dipoles and Diamagnetic (end)
\subsection{Diamagnetism}
Before stepping in diamagnetism, let's perform a cool experiment first. Instead of uniform magnetic field $\vec{B}$, now I use a diverging source and its intensity gradually \textbf{increases}.
\begin{figure}[H]
	\begin{center}
		\includegraphics[width=1\textwidth]{~/PhyAI/chapter6/6_4.png}
	\end{center}
	\caption{A counterclockwise current inside a diverging magnetic field}\label{fig:6_4}
\end{figure}

\begin{figure}[H]
	\begin{center}
		\includegraphics[width=1\textwidth]{~/PhyAI/chapter6/6_5.png}
	\end{center}
	\caption{A clockwise current inside a diverging magnetic field}\label{fig:6_5}
\end{figure}
\begin{figure}[H]
	\begin{center}
		\includegraphics[width=0.95\textwidth]{~/PhyAI/chapter6/6_6.png}
	\end{center}
	\caption{Two loops system tends to be \textbf{repelled away} from the north pole}\label{fig:6_6}
\end{figure}
Inside Bohr's atomic model, each electron has its own unique quantity called \textbf{spin}. Spin determines the direction of an electron's motion (clockwise or counterclockwise). If a material has perfectly matched "loops"\footnote[1]{Chemistry textbooks only go as far as explaining how molecules exhibit (or lack) diamagnetic property by drawing \textbf{molecules orbital diagram}; and now we fully understand why "paired" electrons can cause diamagnetism.} (electron orbits around nucleus), it tends to be \textbf{repelled from stronger magnetic field}; and we call such a material \textbf{diamagnetic}\footnote[2]{Most biological bodies are diamagnetic materials, meaning that we can make them float in mid-air using a sufficiently strong magnetic field: \href{https://www.youtube.com/results?search_query=diamagnetism+frog}{Frog floats in mid-air experiment}.}.
\section{Paramagnetism}\label{sec:paramagnetism} % (fold)
\subsection{The Macroscopic View of Paramagnetism}\label{sub:macroscopic_view_of_paramagnetism} % (fold)
If a material does not have any matched "loops" (or "unpaired" electrons), it tends to be \textbf{attracted from stronger magnetic field}, thanks to the effect of magnetic moment $\vec{\mu}$\footnote[1]{A very well-known example of a paramagnetic material is liquid oxygen, in contrast to liquid nitrogen, a diamagnetic material. You should draw their molecules orbital diagram to see the difference.}.
\begin{figure}[H]
	\begin{center}
		\includegraphics[width=0.75\textwidth]{~/PhyAI/chapter6/6_7.png}
	\end{center}
	\caption{The macroscopic view of paramagnetic material}\label{fig:6_7}
\end{figure}
An external magnetic field $\vec{B}$ immediately magnetizes the material, forcing the random vectors to align; but when the magnetic field is removed, the material returns to its original state. We define the concept of \textbf{magnetization} of a sample, which represents the density of magnetizing dipole moment per unit volume.
\begin{definition}[Magnetization Definition]
	$$ M=\frac{\text{measured magnetic moment }}{V} $$
\end{definition}
In 1985, Pierre Curie discovered experimentally the relationship:
$$ M\propto \frac{B }{T} $$
This proportional relationship seems like very obvious; the higher the temperature, the harder the alignment of magnetic moment vectors, and vice versa.
% subsection Macroscopic view of Paramagnetism (end)
\subsection{The Microscopic View of Paramagnetism}
To examine the paramagnetic material, and explain some surprising phenomenons, we have to equip statistical tool. For simplicity, I restrict our concern to \textbf{two-state paramagnet}. Two-state paramagnet are the materials in which the magnetic moment vector can only point either up or down.
\begin{figure}[H]
	\begin{center}
		\includegraphics[width=0.5\textwidth]{~/PhyAI/chapter6/6_8.png}
	\end{center}
	\caption{Two-state paramagnet}\label{fig:6_8}
\end{figure}
If I denote $N$ as "the number of", the multiplicity and entropy of the system are:
$$ \Omega (N_{\uparrow})=\binom{N}{N_{\uparrow}}\Rightarrow S=k_{B}\ln\Omega=k_{B}\ln\binom{N}{N_{\uparrow}}$$
The internal energy is determined by:
$$ U=\mu B(N_{\downarrow}-N_{\uparrow})=\mu B(N-2N_{\uparrow}) $$
I neglect the constant $B$ in the magnetization formula, and redefine the $M$ quantity as:
$$ M=\mu (N_{\uparrow}-N_{\downarrow})=-\frac{U }{B}$$
Before we look for an analytical solution, let's consider the visual representation for these things. We all know how binomial function work, so I can draw a simple table with $4$ columns:
\begin{figure}[H]
	\begin{center}
		\includegraphics[width=0.5\textwidth]{~/PhyAI/chapter6/6_9.png}
	\end{center}
	\caption{Relationship between $N_{\updownarrow}, S,U,M$}\label{fig:6_9}
\end{figure}
If our two-state paramagnet starts moving from state $N_{\uparrow}=0$ to $N_{\uparrow}=N$, I can plot the graph of $(U,S)$:
\begin{figure}[H]
	\begin{center}
		\includegraphics[width=0.8\textwidth]{~/PhyAI/chapter6/6_10.png}
	\end{center}
	\caption{The graph of $(U,S)$}\label{fig:6_10}
\end{figure}
What can we see here? Recall the definition of temperature (Definition ~\ref{The Third Definition of Temperature}), we determine the temperature of our system at $4$ special labeled points. Wow, this is the first time we observe the \textbf{negative temperature} appears; and seems like these states $U\to 0^+$ and $U\to 0^-$ are very unstable, corresponds to $N_{\upward}\to (N/2)^-$ and $N_{\upward}\to (N/2)^+$. Maybe the "exact state", where $S$ is exactly equal to zero, can not happen.
% section Paramagnetism (end)
% chapter Magnetism of Matters (end) 
\end{document}
